\chapter{Environmental and Waste Management Risk Assessment}
\label{chap:Environmental and Waste Management Risk Assessment}

%---------------------------- SECTION ----------------------------
\section{Why this assessment matters in construction}

Construction is the UK's largest waste-producing sector, responsible for approximately one-third of all waste generated in England. Beyond the volume of waste, construction activities produce a range of environmental pollution risks --- water contamination from surface runoff, soil contamination from fuel spills, air quality impacts from dust and diesel emissions, noise impacts on surrounding communities, and habitat damage from groundwork in ecologically sensitive areas --- that are regulated by environmental legislation with criminal enforcement consequences. The Environment Agency (EA), Natural Resources Wales (NRW), and the Scottish Environment Protection Agency (SEPA) have powers to impose unlimited fines, issue stop notices, and pursue criminal prosecution against contractors and their directors for environmental offences.

The interaction between construction activities and watercourses is a particularly high-risk area. Concrete washout, cement laitance, disturbed clay soils, and hydrocarbon fuel spills all have the potential to reach watercourses through surface water drainage systems, direct overland flow, or groundwater pathways. The pH of concrete washout is approximately 12 to 13; a significant discharge to a watercourse can cause rapid and extensive mortality of aquatic fauna. The EA treats significant pollution incidents as serious criminal matters; the Environment Act 2021 has strengthened enforcement powers and increased financial penalties.

The Duty of Care for waste under the Environmental Protection Act 1990 imposes obligations on every person who produces, holds, or manages waste. Hazardous waste in construction --- contaminated soil, asbestos waste, lead paint, solvent containers, and fluorescent tubes --- must be managed separately from non-hazardous waste, transported by a registered carrier, and disposed of at a licensed facility. The consignment note system for hazardous waste provides the paper trail that enables regulators to trace waste from point of production to point of disposal; failure to maintain accurate consignment notes is an enforcement target.

Beyond compliance, effective environmental risk management on construction sites reduces costs through waste minimisation, avoids programme disruption from pollution investigations and enforcement orders, and protects the project from reputational damage. The assessment must balance the regulatory minimum with the commercial and reputational incentive to manage environmental risk proactively.

\barquote{``A concrete washout that reaches a salmon river is not a regulatory paperwork issue. It is an ecological crime with consequences that extend for years beyond the immediate incident.''}

\example{Site Reality}{A highways contractor is laying a new road junction. The site has a temporary drainage network of interceptor channels draining to a sump at the site boundary before discharging to the adjacent drainage ditch. After three consecutive days of heavy rain, the sump overflows with turbid cement-contaminated water. The overflow enters the drainage ditch, which flows into a designated Site of Special Scientific Interest (SSSI) stream. The EA is notified by a local environmental group; an investigation identifies elevated pH and suspended solids in the stream for 400 metres downstream of the discharge point. The contractor is issued a formal warning and required to fund an ecological assessment of the impact. The construction site pollution prevention plan had identified the sump as adequate for normal rainfall; it had not specified what action should be taken in the event of overflow in sustained heavy rain.}

%---------------------------- SECTION ----------------------------
\section{Legal and professional framework}

The Environmental Protection Act 1990 establishes the Duty of Care for waste and the controls on waste management. The Water Resources Act 1991 and the Water Environment (Water Framework Directive) Regulations 2017 prohibit water pollution and require that activities are managed to prevent contamination of controlled waters (watercourses, groundwater, and coastal waters). The Control of Pollution Act 1974 addresses construction site noise and other nuisances. The Environment Act 2021 has strengthened enforcement powers across all environmental regulation.

The EA's Pollution Prevention Guidelines (PPGs), in particular PPG1 (Understanding Your Responsibilities and Managing Environmental Risk), PPG5 (Works in, Near, or Liable to Effect Watercourses), and PPG22 (Dealing with Spills and Leaks), provide the operational framework for construction pollution prevention. The SEPA and NRW publish equivalent guidance for Scotland and Wales. The Hazardous Waste Regulations 2005 govern the classification, packaging, labelling, consignment noting, and disposal of hazardous construction waste.

\paragraph{Key frameworks relevant to this chapter include: Environmental Protection Act 1990; Water Resources Act 1991; Hazardous Waste Regulations 2005; CDM 2015; EA PPG guidance series.}

\begin{itemize}
  \bitem{Environmental Protection Act 1990 and Duty of Care}{Imposes the Duty of Care on all persons who produce, hold, or manage waste; requires that waste is kept securely; requires that waste is transferred only to an authorised person (registered waste carrier or licensed waste management facility); requires that waste transfer notes are produced for all waste and retained for two years (consignment notes for hazardous waste retained for three years).}
  \bitem{Water Resources Act 1991}{Makes it a criminal offence to cause or knowingly permit the entry of polluting matter into controlled waters; no consent, intention, or negligence is required for the offence to be committed --- causing the pollution is sufficient; strict liability applies.}
  \bitem{Hazardous Waste Regulations 2005}{Require that hazardous construction waste (asbestos, contaminated soils, hydrocarbon-contaminated materials, lead paint, solvent containers, fluorescent lamps) is classified, segregated, labelled, consigned with a consignment note, transported by a registered carrier, and disposed of at a licensed facility.}
  \bitem{CDM 2015}{Requires the Construction Phase Plan to address the environmental management arrangements, including pollution prevention, waste management, and the management of contaminated ground; the pre-construction information must identify any known environmental constraints including watercourse proximity, ground contamination, and ecological designations.}
\end{itemize}

%---------------------------- SECTION ----------------------------
\section{Conducting the assessment using the five-step method}

\subsection{Step 1 --- Identify hazards}
\paragraph{Environmental hazards in construction include: water pollution from cement, concrete, or concrete washout water entering watercourses or drainage systems; hydrocarbon spills from plant refuelling, oil storage, or hydraulic system leaks reaching controlled waters; silt and suspended solids runoff from disturbed ground surfaces reaching watercourses; dust generated by site activities causing air quality nuisance and potential health impacts on receptors adjacent to the site; noise from construction plant and activities causing statutory nuisance for adjoining occupiers; light pollution from site floodlighting causing nuisance to residential neighbours; release of asbestos fibres from poorly managed ACM removal activities; ecology impacts from excavation, vegetation clearance, and dewatering adjacent to habitats and watercourses; and hazardous waste mixed with general waste causing unlawful disposal and potential environmental contamination.}

\subsection{Step 2 --- Decide who or what is affected}
\paragraph{Receptors affected by construction environmental hazards include: watercourses (rivers, streams, drainage ditches) and their associated aquatic and riparian ecology; groundwater bodies used for drinking water abstraction or associated with ecological watercourses; residential occupiers of neighbouring properties affected by noise, dust, light, and vibration; ecological sites including SSSIs, Local Wildlife Sites, and ancient woodland affected by dewatering, drainage changes, or contaminated runoff; the public using adjacent footpaths, roads, or open spaces; and the construction workers themselves, who are environmental law obligors as well as health and safety risk receptors.}

\subsection{Step 3 --- Evaluate likelihood and consequence}
\paragraph{The consequence of a significant pollution incident (contamination of a SSSI watercourse, large hydrocarbon spill) is High in terms of ecological damage, regulatory enforcement, and programme disruption. The likelihood of a pollution incident on an unmanaged construction site adjacent to watercourses is High, given the frequency with which concrete mixing, fuel handling, and disturbed ground runoff occur. Inherent risk is therefore High (16--20/25). A comprehensive site pollution prevention plan with bunded fuel storage, concrete washout facilities, silt fencing, and site drainage management can reduce residual risk to Low to Medium.}

\subsection{Step 4 --- Select and implement controls}
\paragraph{Water pollution controls: install silt fencing and interceptor channels around all disturbed ground areas before any earthworks commence; provide a purpose-built concrete washout facility at the concrete placement area; provide bunded double-skin storage for all fuel; install interceptor drains between the site and watercourses; never discharge concrete washout, cement water, or hydrocarbon-contaminated water to site drains or watercourses. Waste management controls: segregate waste streams at source (general waste; recyclable materials; hazardous waste); engage a registered waste carrier; issue waste transfer notes for all waste collections; issue consignment notes for all hazardous waste; track waste quantities for the site waste management record. Noise and dust controls: implement noise and dust suppression measures referenced in Chapter~\ref{chap:Noise and Vibration Risk Assessment} and Chapter~\ref{chap:Silica, Wood Dust and Respiratory Hazard Risk Assessment} respectively.}

\subsection{Step 5 --- Record and review}
\paragraph{Environmental records on site must include: waste transfer notes for all non-hazardous waste; consignment notes for all hazardous waste; fuel delivery and daily fuel consumption records (for spill investigation purposes); concrete delivery records; and any pollution incident records. The environmental risk assessment must be reviewed when new activities commence, when site drainage arrangements change, and following any pollution incident.}

\example{Five-Step Application}{A demolition contractor is clearing a brownfield site that includes underground fuel storage tanks. Step 1 identifies: fuel-contaminated soil from the tank area; concrete rubble and asbestos-containing materials in the demolition debris; surface water runoff from disturbed ground adjacent to a drainage ditch. Step 2 identifies: local watercourse; residential properties within 50 metres; contractor's employees. Step 3 rates inherent environmental risk as High (16/25). Step 4 implements: fuel-contaminated soils classified as hazardous, excavated separately, tested, and disposed of at a licensed waste facility on consignment notes; asbestos waste removed under licensed conditions before demolition proceeds; silt fence and interceptor drain installed between the site boundary and the drainage ditch before ground is broken; registered waste carrier engaged and waste transfer notes issued. Residual risk rated Low (6/25). Step 5: daily inspection of silt fence integrity and interceptor drain operation.}

%---------------------------- SECTION ----------------------------
\section{Applying the hierarchy of controls}

\paragraph{Environmental risk controls in construction follow the same hierarchy as health and safety controls; elimination of the pollution source or pathway at the design and planning stage is more effective than remediation after a pollution incident.}

\begin{itemize}
  \bitem{Elimination}{Design the construction method to eliminate pollution pathways: specify pre-mixed concrete delivered to site rather than site mixing to eliminate cement mixing washout risk; specify offsite fabrication to reduce the quantity of cutting fluid, oil, and chemical waste generated on site; specify non-contaminating materials where hazardous alternatives could be used.}
  \bitem{Substitution}{Replace hazardous materials with less hazardous equivalents: biodegradable hydraulic oils for plant operating adjacent to watercourses; water-based rather than solvent-based products; non-hazardous concrete release agents.}
  \bitem{Engineering controls}{Bunded double-skin fuel and chemical storage; purpose-built concrete washout facility with sealed settlement system; silt fencing and interceptor channels around all disturbed areas; interceptor drains between the site and watercourses; acoustic barriers and screens for noise-sensitive receptors; dust suppression systems (water bowser, boundary water curtains).}
  \bitem{Administrative controls}{Site Waste Management Plan identifying waste streams, quantities, and disposal routes; waste segregation system with clearly labelled containers; designated waste carrier contact arrangements; daily site environmental inspection; spill response plan with spill kit at fuel and chemical storage areas; noise working hour restrictions; community liaison for noise-sensitive neighbours.}
  \bitem{PPE}{Chemical-resistant gloves and eye protection for fuel and chemical handling; suitable PPE for asbestos and contaminated soil management (as per Chapters~\ref{chap:Asbestos Management Risk Assessment} and \ref{chap:COSHH Assessment for Construction Chemicals}).}
\end{itemize}

%---------------------------- SECTION ----------------------------
\section{Cadence of assessment and review triggers}

\paragraph{Environmental risk assessments must be updated throughout the construction phase as the site develops and the environmental risk profile changes.}

\begin{itemize}
  \bitem{Before any earthworks commence}{The site pollution prevention plan must be in place, silt control measures installed, and fuel and chemical storage arranged before any ground is disturbed.}
  \bitem{Following any weather event that may affect site drainage}{After significant rainfall, the site drainage system must be inspected to confirm that interceptors and silt fences are functioning; any overflow or bypass must be remedied before work resumes.}
  \bitem{When new waste streams are introduced}{When a new activity introduces a new category of waste (hazardous soil, new chemical product, asbestos debris), the waste management arrangements must be reviewed to confirm that the new waste stream is correctly segregated, classified, and consigned.}
  \bitem{Following any pollution incident}{Any spill, release, or water pollution incident must trigger an investigation, immediate regulatory notification where required, and a review of the pollution prevention measures.}
  \bitem{At programme milestones}{The environmental risk assessment must be reviewed at each significant programme milestone, in particular when earthworks are completed and the site drainage pattern changes.}
\end{itemize}

%---------------------------- SECTION ----------------------------
\section{Common failure modes}

\begin{itemize}
  \bitem{Concrete washout discharged to surface water drain}{Concrete truck washout water is discharged onto the site and allowed to run into a surface water gully or drainage ditch without treatment; the alkaline water reaches a controlled water body. This is one of the most common pollution incidents on construction sites.}
  \bitem{Hazardous waste mixed with general waste}{Asbestos debris, solvent containers, and fluorescent tubes are placed in the general waste skip and collected by a general waste contractor; the contractor does not hold a hazardous waste disposal licence; the waste is improperly disposed of.}
  \bitem{Waste transfer notes not issued or retained}{Waste is collected by a skip company and waste transfer notes are not issued; the principal contractor has no evidence of where the waste has gone and cannot produce the notes to the EA on inspection.}
  \bitem{Fuel storage unbunded}{Diesel tanks and oil drums are stored on the site surface without bunding; a fuel leak or accidental spill runs to the site drainage system and reaches a watercourse.}
  \bitem{Silt control measures not installed before earthworks begin}{Earthworks commence before silt fences and interceptor channels are installed; the first significant rainfall event causes turbid runoff to reach the adjacent watercourse.}
  \bitem{Environment Agency notification not made after a pollution incident}{A pollution incident occurs; the contractor attempts to deal with it internally without notifying the EA; the EA discovers the incident independently through downstream sampling; the failure to notify aggravates the regulatory response.}
\end{itemize}

\example{Failure Pattern}{A domestic extension contractor is excavating footings for a rear extension. The excavated clay spoil is stockpiled immediately adjacent to the garden boundary fence. The next day's rainfall washes silt from the spoil heap through a gap in the fence into the adjacent property's garden drain, which connects to the surface water sewer, which in turn connects to a local watercourse. A downstream resident notices discolouration of the water and contacts the EA. The EA inspects the site, identifies the spoil heap and unsupported boundary as the source, and serves a remediation notice on the contractor. The risk assessment for the project did not address the management of excavated spoil or the potential for surface water runoff from spoil stockpiles.}

%---------------------------- CHAPTER SUMMARY ----------------------------
\section{Chapter Summary}

\begin{table}[H]
\centering
\begin{tabularx}{\textwidth}{>{\raggedright\arraybackslash}p{3.2cm}X}
\toprule
\textbf{Assessment Element} & \textbf{Operational Requirement} \\
\midrule
Legal/Professional Basis & Environmental Protection Act 1990; Water Resources Act 1991; Hazardous Waste Regulations 2005; CDM 2015. Strict liability for water pollution; Duty of Care for all waste. \\
Method & Apply the full five-step process and document inherent/residual risk. \\
Controls & Site pollution prevention plan before earthworks; bunded fuel storage; concrete washout facility; silt control measures; hazardous waste segregation and consignment; registered waste carrier. \\
Cadence & Before earthworks; after weather events; when new waste streams introduced; following any pollution incident; at programme milestones. \\
Assurance & Use audit, incident learning, and governance reporting to verify effectiveness. \\
\bottomrule
\end{tabularx}
\end{table}

\begin{itemize}
  \bitem{Key takeaway 1}{Water pollution from construction sites carries strict liability under the Water Resources Act 1991; no intention or negligence needs to be proved; causing the pollution is sufficient. Prevention before the event is the only effective approach.}
  \bitem{Key takeaway 2}{Hazardous waste must be segregated, labelled, consigned on a consignment note, transported by a registered carrier, and disposed of at a licensed facility; mixing hazardous waste with general waste is an offence with unlimited financial penalties.}
  \bitem{Key takeaway 3}{The Site Pollution Prevention Plan must be in place and silt control measures installed before any earthworks commence; retrospective installation after a rainfall event is too late.}
  \bitem{Key takeaway 4}{The Environment Agency must be notified promptly of any significant pollution incident; failure to notify aggravates the regulatory response and demonstrates bad faith that will affect the enforcement outcome.}
\end{itemize}

\barquote{In Chapter~\ref{chap:Demolition Risk Assessment}, we turn from this assessment to the next domain in the Prudentia framework.}
